\documentclass[12pt]{article}
\usepackage[utf8]{inputenc}
\usepackage[T1]{fontenc}
\usepackage[margin=1in]{geometry}
\usepackage{setspace}
\usepackage{xcolor}
\usepackage[colorlinks=true, urlcolor=blue!60!black, linkcolor=blue!60!black]{hyperref}
\setlength{\parindent}{2em}
\setlength{\parskip}{1.25em}
\renewcommand{\baselinestretch}{1.0}
\begin{document}
\begin{center}
{\LARGE \textit{Hicklin v. Orbeck}, 437 U.S.\ 518 (1978)\par}
\vspace{0.8em}
{\large Case Brief --- Four-Page Edition\par}
\end{center}
\vspace{1em}

\section{Memory Jogger}

Under the Article IV Privileges and Immunities Clause, a State may not leverage its ownership of natural resources to force private employers throughout the economic ripple effect of resource development to prefer all state residents over nonresidents when nonresidents are not a ``peculiar source of the evil'' of local unemployment and the preference is not closely tailored (\href{https://supreme.justia.com/cases/federal/us/437/518/}{437 U.S.\ 518}).

\section{Detailed Case Facts}

Citing ``uniquely high unemployment,'' Alaska in 1972 enacted ``Alaska Hire,'' requiring all state oil and gas leases, pipeline permits, and unitization agreements to mandate employing ``qualified Alaska residents'' in preference to nonresidents. Regulations required laying off all nonresidents before any resident in the same craft and administered the preference through ``resident cards.'' The Act covered ``all employment which is a result of'' such leases, reaching contractors, subcontractors, suppliers, and downstream refineries. A ``resident'' needed one continuous year of physical presence and intent to remain.

The Act was not seriously enforced until 1975, at the peak of Trans-Alaska Pipeline construction. In 1976 Labor Commissioner Edmund Orbeck ordered unions to dispatch ``all qualified Alaska residents before any nonresidents.'' Eight individuals, including Sidney Hicklin, most of whom had worked on the pipeline, were barred from pipeline jobs because they could not obtain resident cards.

\section{Procedural History}

The plaintiffs sued in Alaska Superior Court under the Privileges and Immunities and Equal Protection Clauses; the court upheld the Act. The Alaska Supreme Court invalidated the one-year durational residency requirement but, 3--2, upheld the general resident preference, relying on state ownership of oil and gas under \textit{McCready v. Virginia}. The U.S. Supreme Court reversed on June 22, 1978, finding the case not moot as to nonresident appellants.

\section{Judicial Votes}

Reversed, \textbf{9--0}.
\begin{itemize}
\item \textbf{Majority:} Justice Brennan, joined by Chief Justice Burger and Justices Stewart, White, Marshall, Blackmun, Powell, Rehnquist, and Stevens.
\item \textbf{Concurrences and dissents:} None.
\end{itemize}

\section{Holding}

Alaska Hire violates Article IV, \S\ 2. A State may not discriminate against out-of-state citizens in private employment unless there is a ``substantial reason'' (nonresidents are a ``peculiar source of the evil'') and the discrimination bears a ``substantial relationship'' to that evil. State ownership of a resource is a factor, not a categorical exemption, and cannot justify pervasive discrimination imposed on private employers attenuated from the State's proprietary interest, especially for resources bound for interstate commerce. The Court did not reach equal protection and left open whether a preference limited to unemployed residents could survive.

\section{Analysis of Opinions}

\textbf{Majority (Brennan).} The Clause establishes a ``norm of comity'' protecting the right to pursue a common calling in another State. Applying \textit{Toomer v. Witsell}, the record, including a 1971 Federal Field Committee report, showed Alaska's unemployment stemmed from inadequate education, lack of training, and remoteness, particularly among Alaska Natives, not from an influx of nonresidents. The flat preference was grossly overbroad, benefiting a skilled, employed Alaskan as much as an unskilled, unemployed one in job training.

Rejecting the \textit{McCready} defense, Brennan relied on \textit{Baldwin v. Montana Fish and Game Comm'n}: ownership is one factor. The Act swept in private employers with no contract with the State, no work on state land, and no state funding.

Finally, he invoked the ``mutually reinforcing relationship'' between Article IV and the Commerce Clause, both descended from the Articles of Confederation. North Slope oil and gas were destined for interstate commerce and of ``profound national importance,'' and letting a State hoard their employment benefits would invite retaliation and undermine the premise that the States ``must sink or swim together.''

\textbf{Separate opinions.} None; all nine Justices joined in full.

\section{Examples: Future Applications}

\begin{enumerate}
\item \textbf{Same-Side: Wind-Farm Supply Chains.} Wyoming requires wind lessees on state land, their contractors, and suppliers to hire residents first. Land ownership cannot justify coercing downstream private employers, and a flat preference fails tailoring.
\item \textbf{Same-Side: Residency for Licensure.} A State bars nonresidents from practicing law. As in \textit{Supreme Court of New Hampshire v. Piper} (1985), excluding them from a common calling without showing they are a peculiar source of evil is invalid.
\item \textbf{Opposite-Side: Recreational Hunting Fees.} Montana charges nonresidents more for elk licenses. Under \textit{Baldwin}, recreation is not a protected common calling.
\item \textbf{Opposite-Side: State Policymaking Jobs.} A State requires its own troopers and cabinet officials to be residents. Direct employment in sovereign functions implicates the State's political community.
\item \textbf{Fence-Sitter: City Public-Works Quota.} A distressed city requires contractors on city-funded projects to hire 40\% residents. As in \textit{United Building \& Construction Trades Council v. Camden} (1984), the proprietary interest is stronger, but the burden on private employment requires a fact-intensive inquiry.
\end{enumerate}

\section{Critique}

\textbf{Originalist and textualist views.} Some originalist defenders of federalism criticize \textit{Hicklin} and \textit{Piper} for imposing ``unfettered national markets for occupations'' at the expense of state control over local regulation and state resources. Others, such as William Crosskey and John Hart Ely, criticize the inherited ``fundamental rights'' filter from \textit{Corfield v. Coryell} as an atextual gloss, supporting the outcome but questioning malleable balancing.

\textbf{Progressive and doctrinal views.} Commentators praise the decision for preventing beggar-thy-neighbor labor walls, but note that blending Article IV with dormant Commerce Clause cases, which have a market-participant exception Article IV lacks, left uncertainty after \textit{Camden}. Others note that \textit{Hicklin} leaves resource-rich States few tools to channel boom jobs to undertrained local populations such as rural Alaska Natives.

\section{Key Quotations}

\begin{quote}
``In sum, the Act is an attempt to force virtually all businesses that benefit in some way from the economic ripple effect of Alaska's decision to develop its oil and gas resources to bias their employment practices in favor of the State's residents. We believe that Alaska's ownership of the oil and gas that is the subject matter of Alaska Hire simply constitutes insufficient justification for the pervasive discrimination against nonresidents that the Act mandates.''
\end{quote}
\noindent --- Justice Brennan, Opinion of the Court, 437 U.S.\ at 531.

\begin{quote}
``Rather than placing a statute completely beyond the Clause, a State's ownership of the property with which the statute is concerned is a factor---although often the crucial factor---to be considered in evaluating whether the statute's discrimination against noncitizens violates the Clause.''
\end{quote}
\noindent --- Justice Brennan, Opinion of the Court, 437 U.S.\ at 529.

\begin{quote}
``Alaska Hire simply grants all Alaskans, regardless of their employment status, education, or training, a flat employment preference for all jobs covered by the Act. A highly skilled and educated resident who has never been unemployed is entitled to precisely the same preferential treatment as the unskilled, habitually unemployed Arctic Eskimo enrolled in a job training program.''
\end{quote}
\noindent --- Justice Brennan, Opinion of the Court, 437 U.S.\ at 527--28.

\end{document}
